\section{Préparation de la rencontre}

\textbf{Objet :} proposer un cadre fonctionnel à valider avec les personnes demandant la fabrication de l'application.

\textbf{Statut :} document de travail. Les propositions marquées \etatAValider{} ne constituent pas encore des décisions.

\subsection{Objectifs de la rencontre}

La rencontre doit permettre de confirmer :

\begin{itemize}
  \item les rôles des personnes qui créent, instruisent, valident et administrent une DAE ;
  \item les informations visibles ou modifiables selon chaque rôle ;
  \item le cycle complet d'une DAE, y compris les refus, les délais et les relances ;
  \item la politique de conservation et d'archivage ;
  \item la stratégie de reprise des utilisateurs, rôles, services et anciennes demandes ;
  \item le comportement attendu pour créer une nouvelle DAE à partir d'une DAE existante.
\end{itemize}

Les choix techniques déjà présents dans l'application — autorisation serveur, droits au niveau des objets métier, workflow par séquences et séparation des commentaires — servent de base, mais ne remplacent pas la validation métier.

\subsection{Rôles proposés}

Les intitulés suivants correspondent aux rôles actuellement représentés dans l'application. Ils doivent être confirmés avec les utilisateurs :

\begin{table}[htbp]
\centering
\small
\begin{tabularx}{\textwidth}{>{\bfseries}p{0.27\textwidth}Y}
\toprule
\rowcolor{bleuDAE}
\tableauTitre{Rôle proposé} & \tableauTitre{Responsabilité principale} \\
\midrule
Administrateur & Administrer les utilisateurs, les référentiels, les droits et, si nécessaire, intervenir sur une DAE visible. \\
Demandeur & Créer une DAE, la compléter et la lancer dans le circuit. \\
Valideur de pôle & Examiner la cohérence opérationnelle de la demande. \\
Valideur RRH & Examiner les aspects ressources humaines. \\
Valideur DRH & Effectuer la validation RH de niveau supérieur. \\
Valideur contrôle de gestion & Vérifier le financement et les éléments budgétaires. \\
Valideur direction & Rendre la décision finale lorsque le circuit le prévoit. \\
\bottomrule
\end{tabularx}
\caption{Rôles proposés et responsabilités principales.}
\label{tab:roles-proposes}
\end{table}

Un même utilisateur peut être rattaché à plusieurs services. Son rôle principal et son périmètre de services doivent être distingués : le rôle détermine ce qu'il peut faire, le service détermine quelles données il peut voir.

\subsection{Proposition de droits}

La proposition détaillée figure dans l'annexe. Le principe recommandé est le suivant :

\begin{itemize}
  \item le demandeur peut modifier sa DAE tant qu'elle n'est pas lancée ;
  \item les valideurs consultent les données nécessaires à leur décision et ne modifient pas les données du demandeur, sauf règle explicitement validée ;
  \item l'administrateur peut intervenir selon sa visibilité sur la DAE, mais toute intervention est tracée ;
  \item les informations d'audit, les votes et les étapes du workflow ne sont jamais modifiables comme de simples champs de formulaire ;
  \item le contrôle serveur est obligatoire, même lorsqu'un champ est masqué ou désactivé dans l'interface.
\end{itemize}

\begin{tcolorbox}[enhanced,colback=vertClairDAE,colframe=vertDAE,boxrule=0.6pt,arc=1.5mm]
\recommande{} : présenter les droits pendant la réunion avec des noms d'objets métier compréhensibles, et réserver les identifiants techniques à la documentation d'implémentation.
\end{tcolorbox}

\subsection{Cycle proposé d'une DAE}

\begin{table}[htbp]
\centering
\small
\begin{tabularx}{0.94\textwidth}{>{\bfseries}p{0.24\textwidth}Y}
\toprule
\rowcolor{bleuDAE}
\tableauTitre{État} & \tableauTitre{Évolution proposée} \\
\midrule
Brouillon & Modification par le demandeur avant lancement. \\
Lancement du circuit & Copie de la configuration applicable au service. \\
Validation séquence par séquence & Chaque rôle examine et décide selon son étape. \\
Refusée & Un refus motivé clôt immédiatement le circuit. \\
Acceptée & Toutes les séquences prévues ont été acceptées. \\
Archivée & La DAE reste consultable en lecture seule selon la règle de conservation. \\
\bottomrule
\end{tabularx}
\caption{Cycle proposé et évolution des états d'une DAE.}
\label{tab:cycle-synthetique}
\end{table}

\subsubsection{Règles proposées}

\begin{enumerate}
  \item Le demandeur crée et enregistre la DAE comme brouillon.
  \item Le demandeur lance le circuit ; un administrateur peut le faire à sa place s'il voit la DAE.
  \item Au lancement, les étapes et rôles du service sont copiés dans l'exécution du circuit. Une modification ultérieure du paramétrage ne change pas le circuit déjà lancé.
  \item Chaque rôle de la séquence doit accepter avant l'ouverture de la séquence suivante.
  \item Un refus motivé clôt le circuit.
  \item Les relances sont envoyées lorsque le délai configuré est dépassé.
  \item Un administrateur habilité peut suppléer un rôle après expiration de son délai.
  \item Les commentaires généraux de la DAE restent séparés des commentaires de décision.
\end{enumerate}

\subsection{Visibilité et archivage — décision attendue}

\begin{table}[htbp]
\centering
\small
\begin{tabularx}{\textwidth}{>{\bfseries}p{0.25\textwidth}Yp{0.18\textwidth}}
\toprule
\rowcolor{bleuDAE}
\tableauTitre{Option} & \tableauTitre{Effet} & \tableauTitre{Avis} \\
\midrule
Conservation visible en permanence & Toutes les DAE restent dans la liste, avec filtres par état. & Simple, mais la liste peut devenir difficile à exploiter. \\
Archivage logique configurable & Après une durée validée, la DAE reste consultable en lecture seule dans les archives. & \recommande{} \\
\bottomrule
\end{tabularx}
\caption{Options de visibilité et d'archivage des DAE.}
\label{tab:options-archivage}
\end{table}

\begin{tcolorbox}[enhanced,colback=orangeClairDAE,colframe=orangeDAE,boxrule=0.6pt,arc=1.5mm]
\surligneOrange{Questions à trancher :} après combien de temps une DAE acceptée, refusée ou abandonnée est-elle archivée ? Les DAE refusées suivent-elles la même durée ? Qui peut consulter les archives ? Une DAE archivée peut-elle être désarchivée, ou seulement dupliquée ? Existe-t-il une durée légale ou métier de conservation ?
\end{tcolorbox}

La suppression physique est déconseillée : elle ferait perdre l'historique et rendrait les contrôles ultérieurs plus difficiles.

\subsection{Reprise des anciennes données}

Le dépôt ne contient pas le schéma de l'ancienne application. La migration doit donc commencer par une phase d'inventaire et de rapprochement, et non par un mapping définitif supposé.

La démarche proposée est :

\begin{enumerate}
  \item obtenir un export ou un accès en lecture à l'ancienne base ;
  \item inventorier les utilisateurs, rôles, services, rattachements, DAE et historiques de validation ;
  \item dédoublonner les utilisateurs par courriel, identifiant historique puis nom et prénom ;
  \item rapprocher les anciens rôles des rôles cibles ;
  \item reprendre l'ensemble des rattachements d'un utilisateur à ses services ;
  \item produire un rapport des correspondances incertaines ou impossibles ;
  \item importer d'abord dans un environnement de test ;
  \item contrôler les volumes, les doublons, les utilisateurs sans rôle et les services disparus ;
  \item valider le résultat avant la reprise en production.
\end{enumerate}

Les mots de passe, comptes inactifs et données personnelles particulières feront l'objet d'une décision spécifique. Les secrets ne doivent pas être exportés dans les fichiers de migration ou de documentation.

\subsection{Duplication d'une DAE existante}

La fonction proposée est : \surligneBleu{« Créer une nouvelle DAE à partir de cette DAE »}.

Elle crée une nouvelle demande indépendante, dans l'état brouillon, appartenant à l'utilisateur courant. Elle ne relance jamais automatiquement l'ancien circuit et ne copie aucun vote.

\textbf{À copier, sous réserve que les références soient encore actives :}

\begin{itemize}
  \item service concerné ; type de contrat ; motif de recrutement ; motif de remplacement ; financement ;
  \item intitulé du poste ; catégorie professionnelle ; temps de travail et répartition ;
  \item rémunération et périodicité ; qualification ; indicateurs métier ;
  \item candidat ou collaborateur, uniquement si cette information reste pertinente.
\end{itemize}

\textbf{À réinitialiser ou laisser vide :}

\begin{itemize}
  \item identifiant de la DAE et demandeur, remplacé par l'utilisateur courant ;
  \item dates de création et de modification ; état de validation ; étapes, visas et votes ;
  \item commentaires généraux et commentaires de décision ; date d'archivage ; date de suppression ;
  \item résultat du précédent circuit.
\end{itemize}

Si une référence copiée est inactive ou supprimée, elle doit être laissée vide et signalée au demandeur pour ressaisie.

\subsection{Décisions à obtenir pendant la rencontre}

\begin{itemize}
  \item confirmer les rôles et leurs intitulés métier ;
  \item valider la matrice des droits objet par objet ;
  \item confirmer qui peut lancer une DAE au nom du demandeur ;
  \item confirmer les délais, relances et règles de suppléance ;
  \item choisir la politique et la durée d'archivage ;
  \item confirmer les données à reprendre dans les anciennes DAE ;
  \item confirmer les données à recopier ou remettre à zéro lors d'une duplication ;
  \item identifier les règles de conservation imposées par le métier ou la réglementation.
\end{itemize}

\subsection{Références dans l'application}

\begin{itemize}
  \item ADR 0042 — autorisation au niveau du champ ;
  \item ADR 0048 — cinq zones fonctionnelles de la DAE ;
  \item ADR 0050 — périmètres hiérarchiques des utilisateurs ;
  \item ADR 0055 — workflow de validation ;
  \item ADR 0056 — remise à zéro exceptionnelle d'un cycle ;
  \item ADR 0062 — export PDF sécurisé des DAE.
\end{itemize}
